\documentclass[10pt,a4paper]{amsart}
\usepackage[margin=19mm]{geometry}
\usepackage{amsmath,amssymb,mathtools}
\usepackage[scr=rsfs,cal=euler]{mathalfa}
\usepackage{xfrac}
\usepackage[unicode,hidelinks,bookmarks=false]{hyperref}
\newcommand{\ie}{i.e.\ }
\newcommand{\cf}{cf.\ }
\newcommand{\resp}{respectively\ }
\newcommand{\Subsection}{\subsection}
\providecommand{\nobreakdash}{\nobreak\hbox{-}}
\setlength{\emergencystretch}{2em}
\multlinegap0pt
\usepackage{amssymb}
\usepackage{xcolor}
\usepackage{tikz}
\newtheorem{thm}{Theorem}
\newtheorem{lemma}{Lemma}
\newtheorem{lem}{Lemma}
\title{On separated families of Anosov representations}
\author{Joaqu\'in Lejtreger, Joaqu\'in Lema}
\date{Source: arXiv:2604.18994v2 (CC BY 4.0)}
\begin{document}
\maketitle

\noindent\emph{Source and licence.} On separated families of Anosov representations. Version arXiv:2604.18994v2 (CC BY 4.0), distributed by arXiv under the Creative Commons Attribution 4.0 International licence, http://creativecommons.org/licenses/by/4.0/, as recorded in the arXiv licence metadata for this exact version and shown on the abstract page https://arxiv.org/abs/2604.18994v2. Full licence notice: \texttt{licenses/arxiv-2604.18994-CC-BY-4.0.txt}.\par\smallskip

\noindent\textit{Editorial scope.} Counted unit: the article's thm thm:sl3 (source0.tex lines 1418--1431). The article's complete original proof of it is delivered with it (source0.tex lines 1433--1464, one \texttt{proof} environment of 32 lines). No mathematics is written, repaired or re-derived: the statement and the proof are the article's own lines, in the article's own notation, and no retained line is substituted or re-flowed.  Retained uncounted as the source-local context the proof needs: lines 956--965; lines 1230--1248; lines 1507--1510.  Nothing was omitted from any retained span, so the package carries no omission record.  No retained line contains a citation, so the wrapper carries no bibliography and no bibliography entry.  No retained line contains a figure environment, an image-inclusion command or drawing code, so no image is delivered and the package declares no asset dependency.  Editorial formatting only, in addition: an AMS \texttt{amsart} preamble that restates the article's own declarations and macros used by the retained lines, namely the environments \texttt{thm}, \texttt{lemma}, \texttt{lem} on separate counters, together with the standard packages those lines need.  The retained environments print in this document as Theorem 1, Lemma 1, Lemma 1 in order of appearance, whereas the article prints the same items with the numbering of its own full document.  The complete native source of the article as distributed in the arXiv e-print for this exact version is bundled unchanged as \texttt{sources/198837/source0.tex}.
\begin{thm}
\label{thm:general+}
Let $\Theta \subseteq \Pi$ with, $\varphi \in \mathfrak{a}_\Theta^*$, and $(\rho_n) \in Hom (\Gamma,G)^\mathbb{N}$ be a $\Theta-$strongly separated sequence of representations with respect to a strong Markov coding $\mathcal{G}$ such that $\frac{\min_{g \in S} \varphi (\kappa (\rho_n (g)))}{ -\log \sin \varepsilon_n} \to \infty,$ with $(\varepsilon_n)$ a sequence as above. 
Then
$$\lim_{n \to \infty} \frac{h_{\rho_n} (\varphi)}{h_{\rho_n}^\kappa (\varphi)} = 1.$$
Moreover, assuming $\iota(\Theta) = \Theta$ and that $0 \neq \inf_n h_{\rho_n}^\kappa \varphi (\kappa (\rho_n (g)))$ for any $g \in S$, then
$$\sup_n d_{Th}^\varphi (\rho_n, \rho_0) <\infty,$$
and similarly, if $\sup_n h_{\rho_n}^\kappa \varphi (\kappa (\rho_n (g))) <\infty$ for every $g \in S$, then
$$\sup_n d_{Th}^\varphi (\rho_0,\rho_n) < \infty.$$
\end{thm}

\begin{lemma}
\label{lemma:eigenvalues}
 Let $\mathfrak{a} \subset \mathfrak{sl}_3 (\mathbb{R})$ be the standard Cartan subspace, and $\lambda : SL_3 (\mathbb{R}) \rightarrow \mathfrak{a}^+$ the Jordan projection.
 Then given $g = E(Z_k,W_k) T(X_k) \ldots E(Z_1,W_1) T(X_1)$, such that $\prod_{i=1}^l Z_i^{-1} X_i < 1$ and $\prod_{i=1}^l W_i > 1$, then we get that
 $$\alpha_1 (\lambda (g)) = \sum_{i=1}^l \log W_i, \: \alpha_2 (\lambda (g)) = \sum_{i=1}^l \log Z_i - \sum_{i=1}^l \log X_i .$$
 Moreover, if $W_i^n,Z_i^n$, and $X_i^n$ are sequences verifying the inequalities above, and such that $Z_i^n, W_i^n \to \infty$, $(Z_i^n)^{-1} X_i^n \to 0$, then the stable, central and unstable direction of $g_n = E(Z_k^n,W_k^n) T(X_k^n) \ldots E(Z_1^n,W_1^n) T(X_1^n)$ converge to $e_1$, $e_2$ and $e_3$ respectively. 
\begin{proof}
 The first part of the proof is a straightforward computation.
 As for the second part, notice that under these conditions, the first row of the blocks $E(Z_i^n, W_i^n) T(X_i^n)$ goes to zero, when $n \to \infty$.
 Multiplying such upper triangular matrices, we get something of the form
   $$g_n = 
   \begin{pmatrix}
 o(1) & o(1) & o(\prod_{i=1}^l W_i^n) \\
   0 & 1 & o(\prod_{i=1}^l W_i^n)\\
   0 & 0 & \prod_{i=1}^l W_i^n
   \end{pmatrix}.$$
 The stable, central, and unstable directions of these matrices are approximately $e_1$, $e_2$, and $e_3$, respectively.
\end{proof}
\end{lemma}

\begin{lem}
\label{lemma:locallyconstant}
Let $\phi \in C^{\alpha}(X_A)$ be a locally constant function. We define the \emph{weighted matrix} $A_\phi$ as the matrix obtained by replacing $a_{ij}$ in $A$ by $e^{\phi(i)}$. Then the spectral radius of $A_\phi$ is $e^{P(\phi)}$.
\end{lem}

\begin{thm}
  \label{thm:sl3}
 Let $\varphi \in \mathfrak{a}^*$ and $(\vec{Z_t})_{t \in \mathbb{R}_{\geq 0}}, (\vec{W_t})_{t \in \mathbb{R}_{\geq 0}}$ be a path of vectors such that $\vec{Z_t},\vec{W_t} \to \infty$.
 Then, given $(\rho_t) \in \mathbb{RP}(X_1, X_2)^{\mathbb{R}_{\geq 0}}$ a path of holonomies of convex projective structures with edge invariants $\vec{W_t},\vec{Z_t}$, and assuming $\varphi$ is positive on their limit cone, we get that
   \begin{enumerate}
      \item Let $v_i (t) = (\log (W^i_t Z^{i-1}_t)) \omega_1 + \log (Z^i_t W^{i-1}_t) \omega_2$, $w_i (t) = (\log (Z^i_t W^{i-1}_t)) \omega_1 + (\log (W^i_t Z^{i-1}_t)) \omega_2$, we get that
         $$\lim_{t \to \infty }\frac{h_{\rho_t} (\varphi)}{s(t)} =1,$$
 where $s(t)$ is a positive solution to the equation
         $$\prod_{i=1}^3 (1-e^{-s (t) \varphi (v_i (t))}) + \prod_{i=1}^3 (1-e^{-s (t) \varphi (w_i (t))}) + \sum_{i=1}^3 e^{-s (\varphi (v_i + w_i))} =1.$$
      \item Let $\rho_t \in \mathbb{RP} (X_1, X_2)$, and $\rho_t' \in \mathbb{RP} (X_1', X_2')$ be a path of representations in the constant triple ratio slice, and defined by the same family of cross ratios given by $\vec{Z_t},\vec{W_t}$, then
      $$d_{Th}^\varphi (\rho_t,\rho_t') \to_{t \to \infty} 0.$$
      \item If $\max \{\varphi (v_i (t)),\varphi (w_i (t))\}_{i=1,2,3}/\min \{\varphi (v_i (t)),\varphi (w_i (t))\}_{i=1,2,3}$ is bounded above and below for every $t$, then the Thurston asymmetric metric defined by $\varphi$ is finite.
   \end{enumerate}
\end{thm}

\begin{proof}
\begin{enumerate}
    \item 
 Recall that the automaton defined by the symmetric set of generators $a^{\pm 1}, b^{\pm 1}, c^{\pm 1}$ is conjugated to the shift of finite type for the matrix $A$. 

 By Theorem \ref{thm:general+}, we have that $$\lim_{t \to \infty} \frac{h_{\rho_t}(\varphi)}{h_{\rho_t}^\kappa(\varphi)} = 1,$$
 where $h_{\rho_t}^\kappa(\varphi)$ is the only solution to $P(-sC_\varphi) = 0$, for $C$ the locally constant potential given by the Cartan projection of the first letter of the coding. 
    
 Notice that under the strong separation conditions, $\varphi (\kappa (\rho_n (g)))$ is asymptotic to $\varphi (\lambda (\rho_n (g)))$.
 Particularly, we get that $\lim_{t \to \infty} h_{\rho_t}^\kappa/h_{\rho_t}^\lambda = 1$.
 Using Lemma \ref{lemma:locallyconstant}, we get that $h_{\rho_t}^\lambda$ is the unique number for which the matrix:

 $$\mathcal{L}_s = 
 \begin{pmatrix}
 e^{- s \varphi (w_1 (t))} & 0 & e^{- s \varphi (w_1 (t))} & 0 & e^{- s \varphi (w_1 (t))} & e^{- s \varphi (w_1 (t))}\\
 e^{-s \varphi (w_2 (t))} & e^{-s \varphi (w_2 (t))} & 0 & e^{-s \varphi (w_2 (t))} & 0 & e^{-s \varphi (w_2 (t))}\\
   0 & e^{-s \varphi (w_3 (t))} & e^{-s \varphi (w_3 (t))} & e^{-s \varphi (w_3 (t))} & e^{-s \varphi (w_3 (t))} & 0\\
   0 & e^{-s \varphi (v_1 (t))} & e^{-s \varphi (v_1 (t))} & e^{-s \varphi (v_1 (t))} & e^{-s \varphi (v_1 (t))} & 0\\
 e^{-s \varphi (v_2 (t))} & 0 & e^{-s \varphi (v_2 (t))} & 0 & e^{-s \varphi (v_2 (t))} & e^{-s \varphi (v_2 (t))}\\
 e^{-s \varphi (v_3 (t))} & e^{-s \varphi (v_3 (t))}  & 0 & e^{-s \varphi (v_3 (t))}  & 0 & e^{-s \varphi (v_3 (t))} .
 \end{pmatrix}
 $$
 An algebraic calculation (maybe done by computer) produces an expression for the characteristic polynomial of this matrix.
 Assuming one is a solution, one gets the equation in the first item.
 
      \item Again, since $\varphi (\kappa (\rho_n (g)))$ is asymptotic to $\varphi (\lambda (\rho_n (g)))$, it is enough to estimate the former.
 On the other hand, Lemma \ref{lemma:eigenvalues} shows that the asymptotics for the Jordan projection of $a,b,c$ at constant triple ratio are independent of $(X_1, X_2)$.
 Particularly, the asymptotics for the length function are independent of which slice you are in.
      
   \item This is a direct consequence of the second part of Theorem \ref{thm:general+}.
\end{enumerate}
\end{proof}

\end{document}
